\chapter{Commenti Finali}

\section{Autovalutazione e lavori futuri}

\subsection{Massimiliano Ria}

Nutro un forte interesse per la programmazione ad alto livello e in particolare per il paradigma object-oriented, il che ha reso particolarmente elettrizzante l'idea di approcciarmi
a un progetto interamente sviluppato da me e dai miei colleghi, in modo da poter mettere in pratica gli argomenti affrontati durante il corso in contesti più circoscritti. 

La prima difficoltà incontrata è stata senza dubbio il lavoro di gruppo: data infatti la mia personalità intraprendente e solitaria ho trovato moderatamente ostico dovermi
confrontare con i miei collaboratori prima di adottare una determinata soluzione che mi sembrava potesse funzionare. Tuttavia ritengo che questo fosse proprio uno degli obiettivi principali 
del corso: imparare, con l'esperienza, ad acquisire maggiore confidenza con il lavoro di squadra, una realtà con cui sarà necessario confrontarsi nel mondo del lavoro.

All'interno del progetto mi sono occupato di diversi aspetti del sistema: ho contribuito alla struttura MVC, implementando le componenti necessarie alla comunicazione tra le varie
parti; ho sviluppato l'interfaccia grafica, dalle sue basi fino ai dettagli più avanzati; infine, lato Model, ho progettato il sistema di intelligenza artificiale, dal flusso di comunicazione
ai veri e propri algoritmi di ricerca.

Mi sento di poter dire che nonostante le difficoltà incontrate, i lunghi tempi richiesti e le situazioni di stallo che di certo non sono mancate,
la mera scrittura di codice è stata decisamente divertente, la progressione del sistema altamente soddisfacente e le discussioni avute con i compagni di viaggio senza dubbio uniche.

Mi sento dunque soddisfatto del lavoro svolto: ho cercato di rispettare i pattern e i paradigmi della programmazione a oggetti raccomandati durante il corso, ho sfruttato i vantaggi
della programmazione funzionale quando necessario e sono riuscito a collaborare con gli altri componenti del gruppo facendo qualche passo indietro e mettendomi più spesso in discussione.

In futuro sono certo che avrò nuove occasioni per migliorare la qualità del mio codice e imparare altri linguaggi orientati agli oggetti, che mi risulteranno sicuramente più accessibili
grazie all'esperienza già maturata con Java.

\subsection{Simone Marsili}

Questo è stato sicuramente il progetto più complesso che io abbia realizzato fino a questo momento.

La programmazione object-oriented mi ha affascinato fin da subito e ha scaturito in me la voglia di mettermi alla prova con qualcosa di nuovo.
Il nuovo ha necessariamente portato con sé delle difficoltà, le quali in un primo momento sembravano insormontabili: per esempio l'idea di sviluppare l'architettura del progetto e partire dalla relazione
prima ancora di mettere mano al codice, oppure riuscire a collaborare parallelamente con i compagni usando lo strumento spaventoso noto come Git.
Non sono mancati quindi i momenti di sconforto, soprattutto nella fase iniziale, quando la sintonia con i membri del gruppo era veramente debole e qualsiasi scelta architetturale e progettuale potenzialmente errata
sembrava potesse portare ad un inevitabile fallimento.

Il desiderio di coesione con il team di lavoro, l'interesse verso il videogioco scelto, la voglia di imparare e le grandi soddisfazioni date da piccoli progressi mi hanno dato l'impulso necessario ad andare avanti e a
mettere tutto me stesso nel progetto.

La scrittura in sé del codice non ha rappresentato invece un grosso problema per me, anzi è stata molto piacevole e stimolante.
Mi sono occupato principalmente dello sviluppo degli elementi base del modello, assicurandomi di renderli solide fondamenta a supporto delle altre componenti sviluppate dai miei compagni.
Mi sono anche cimentato nel dialogo MVC, sviluppando in particolare la comunicazione Model -> View e cercando di adottare la scelta progettuale più adatta al contesto del videogioco.
Nel complesso sono abbastanza soddisfatto di come io abbia progettato e sviluppato il codice: ho cercato di seguire in ogni momento ciò che ho imparato durante le lezioni del corso e credo di aver reso le mie parti di progetto
quanto più stabili, pulite, chiare e soprattutto estendibili verso futuri aggiornamenti.

Sono ancor più felice e soddisfatto dei rapporti che ho costruito con gli altri membri del team: abbiamo legato molto in questi mesi e l'aver affrontato le difficoltà ha accentuato la coesione del gruppo,
all'interno del quale ho cercato di mantenere spesso un occhio cinico e critico utile allo sviluppo.

In futuro, mi piacerebbe portare avanti il progetto assieme ai miei compagni, per esempio inserendo le "Spell" all'interno del gioco, e sono sicuro che sarà una challenge facilmente affrontabile grazie alle competenze acquisite e al percorso svolto fin qui.

\subsection{Francesco Dama}

Programmazione ad oggetti è stata sicuramente la materia più interessante che ho incontrato fin'ora nel corso di studi, è stato un nuovo
linguaggio di programmazione da imparare da zero, sicuramente una sfida importante che ho accettato volentieri.
Durante lo svolgimento del corso, principalmente nei laboratori, ci è stato insegnato come lavorare con Git nel singolare, non in gruppo,
quindi questa è stata sicuramente una delle difficoltà che abbiamo affrontato nel corso del progetto.
Essendo in 4 a discutere sul da farsi, la fase che è stata più difficile secondo me è stata la prima, ovvero la fase di discussione del dominio
e del modello, quando ancora non avevamo un'idea materiale e non pensavamo minimamente al codice da scrivere; Avevamo tutti idee contrastanti
che avrebbero portato a dei vicoli cechi.
Il tutto si è risolto quando abbiamo cominciato effettivamente a scrivere codice, codice sul quale non avevamo dubbi, ovvero le fondamenta
del progetto, come i concetti di Deck e Carte.
Quando ho iniziato a scrivere anch'io ho notato che è stato molto più semplice di quanto mi aspettassi, avevo le idee chiare in testa e sono
riuscito a trasformarle correttamente in codice, senza alcuna grande problematica.

Sono molto soddisfatto di ciò che ho scritto, ma avrei preferito mettermi in gioco ancora di più, perchè penso di aver scritto/ideato
poche cose rispetto ai miei compagni, sicuramente quando aggiorneremo il nostro progetto con i concetti opzionali, mi metterò molto di più
in gioco.
Una cosa che mi sento quindi di criticare è la suddivisione del lavoro, apparentemente errata dato che vi è stato qualcuno che ha scritto
molto di più di altri, quando avrei preferito affrontare argomenti più difficili e interessanti, che non ho particolarmente incontrato nella 
parte da me affrontata.

Tutto sommato questo progetto mi ha sicuramente aperto la mente al lavoro di squadra, sia in presenza che da remoto tramite Git, ed è stata
una nuova esperienza che non avevo mai affrontato precedentemente, dove all'inizio sembrava impossibile realizzarla ma poi si è rivelata 
affascinante, grazie sia agli insegnamenti dei docenti sia al linguaggio in sè.

\subsection{Lorenzo Mercuriali}

Come membro del gruppo che ha proposto l'idea di realizzare un progetto ispirato al gioco di carte HearthStone, sono molto contento dello
sviluppo del lavoro svolto, che ha portato il gioco allo stato attuale.

Tuttavia, il progetto ha comportato diverse difficoltà e complessità, insieme a momenti in cui non sapevamo come
muoverci correttamente per proseguire il nostro lavoro, specialmente nelle fasi iniziali; più volte le nostre idee si sono scontrate
e, attraversando momenti critici, il rapporto all'interno del gruppo si è rafforzato notevolmente.

Non avere ancora superato con successo l'esame pratico in laboratorio ha messo in luce alcune problematiche durante la scrittura del codice, ma,
grazie ai miei colleghi, che sono stati sempre disponibili e pronti a darmi consigli, sono riuscito a superare ostacoli che pensavo fossero
invalicabili. Sono soddisfatto di ciò che ho realizzato; tuttavia, mi sarebbe piaciuto mettermi alla prova in maniera ancora più significativa,
poiché ritengo che il mio contributo, specialmente nelle fasi iniziali, non sia stato ottimale e sia diventato più concreto solo dopo che il progetto aveva compiuto
qualche passo avanti.

Infine, ritengo che HearthCode abbia migliorato le mie abilità nella programmazione a oggetti, anche in vista di affrontare nuovamente
l'esame. Non vedo l'ora di poter inserire un giorno, insieme ai miei colleghi, le "Spell" nel gioco di carte, rendendolo ancora
più competitivo, imprevedibile e divertente.
